\documentclass[11pt,a4paper]{article}
\usepackage[utf8]{inputenc}
\usepackage[T1]{fontenc}
\usepackage[brazilian]{babel}
\usepackage[margin=2.5cm]{geometry}
\usepackage{graphicx}
\usepackage{booktabs}
\usepackage{array}
\usepackage{xcolor}
\usepackage{hyperref}
\usepackage{caption}
\usepackage{amsmath}
\usepackage{longtable}
\usepackage{float}

\hypersetup{
  colorlinks=true,
  linkcolor=blue!50!black,
  urlcolor=blue!50!black,
  citecolor=blue!50!black
}

\title{Detecção de Anomalias em Cabiúnas via AutoML\\
\large Experimento EXP5 --- Grupo TC382\_03\_A: histórico, modelo final,\\
antecedência de detecção e falsos positivos}
\author{Pipeline \texttt{cnn1d\_ae} / AutoML --- projeto \texttt{TesteMLCab}}
\date{Agosto de 2026}

\begin{document}
\maketitle

\begin{abstract}
Este relatório documenta a série de experimentos AutoML (\texttt{EXP5}) realizada
sobre o sensor \texttt{TC382\_03\_A} do equipamento Cabiúnas, incluindo o processo
de busca de modelo (\texttt{dense}/\texttt{ocsvm}/\texttt{iforest}), a validação
out-of-sample, a checagem de variância de semente, e uma análise detalhada de
\emph{antecedência de detecção} (quanto tempo antes do alarme real o modelo
sinaliza anomalia) e de \emph{falsos positivos}. O candidato final validado é um
\texttt{IsolationForest} univariado sobre \texttt{TC382\_03\_A}, com taxa de
acerto de 65\% sobre alarmes reais fora da amostra, dos quais 45\% detectados
com antecedência real (mediana de 13,4h antes do alarme), variância de semente
nula, e aproximadamente 1 episódio de falso alerta a cada 3 dias.
\end{abstract}

\tableofcontents

% =====================================================================
\section{Objetivo}
% =====================================================================

Avaliar, dentro da pipeline própria de AutoML (\texttt{src/cnn1d\_ae/automl\_pipeline.py},
adaptada da pipeline \texttt{transpetro\_modelos} da Lara), qual combinação de
modelo/sensor(es)/hiperparâmetros melhor detecta os alarmes reais do
Cabiúnas para o sensor \texttt{TC382\_03\_A} (temperatura), e caracterizar
com que antecedência essa detecção acontece e a que custo em falsos positivos.

% =====================================================================
\section{Dados}
\label{sec:dados}
% =====================================================================

\begin{itemize}
  \item \textbf{Fonte:} \texttt{serie\_consolidada\_2025\_interpolated\_antigo.csv}
        (dataset ClearML \texttt{Cabiunas 2025}, versão \texttt{424e5b58...}),
        resolução de 30s.
  \item \textbf{Período disponível:} 2025-01-01 a 2025-10-31 (10 meses,
        $\sim$875 mil pontos).
  \item \textbf{Sensor analisado:} \texttt{TC382\_03\_A} (univariado nesta
        etapa final; o histórico também cobre um grupo combinado com
        \texttt{T5\_AVG\_A}).
  \item \textbf{Sensor de referência operacional:} \texttt{NGP\_A}, usado para
        derivar o estado operacional (\texttt{on} / \texttt{off\_longo} /
        \texttt{off\_curto} / \texttt{transiente}).
  \item \textbf{Alarmes:} \texttt{alarmes\_record\_2025\_tags\_modelo.csv}.

  \textbf{Achado de qualidade de dado:} o arquivo de alarmes tem
  \emph{cada linha duplicada exatamente 2×} (562 linhas $\to$ 281 alarmes
  reais únicos). Isso não altera nenhuma taxa (\texttt{hit\_rate},
  \texttt{composite\_score}) reportada nas seções~\ref{sec:historico}--\ref{sec:modelo},
  pois a duplicação é uniforme e se cancela na razão -- mas os contadores
  absolutos de \texttt{n\_alarms} gerados pela pipeline até este relatório
  estavam em dobro do real. A análise de antecedência da
  Seção~\ref{sec:antecedencia} já usa os valores deduplicados.
\end{itemize}

% =====================================================================
\section{Metodologia}
% =====================================================================

\subsection{Pipeline AutoML}

Diferente do autoencoder CNN-1D sequencial (janelado) usado no restante do
projeto, a pipeline AutoML opera \textbf{ponto a ponto}: cada instante de
tempo é uma amostra de treino independente, sem janelamento temporal
(\texttt{seq\_len}). Três famílias de modelo são testadas:

\begin{itemize}
  \item \textbf{dense}: autoencoder denso (Keras), reconstrói o vetor de
        features do instante; erro = MSE de reconstrução.
  \item \textbf{ocsvm}: One-Class SVM (kernel RBF); erro = $-$decision\_function.
  \item \textbf{iforest}: Isolation Forest; erro = $-$score\_samples (quanto
        mais isolado, mais anômalo).
\end{itemize}

Para cada modelo, um \emph{grid search} varre percentis de threshold
($\{95, 97, 99, 99.5\}$, calculados sobre o erro de reconstrução do
conjunto de treino) e \texttt{debounce} ($\{1, 3, 6, 12\}$ pontos
consecutivos exigidos), rankeando todos os \emph{trials} por
\texttt{composite\_score}.

\subsection{Métrica de avaliação}

O \texttt{composite\_score} balanceia detecção de alarme real contra falso
alerta:
\[
\text{balanced\_score} = \text{detection\_rate} - \text{fp\_penalty} \cdot \text{normal\_alert\_rate}^2
\]
com penalidade adicional se \texttt{detection\_rate} cair abaixo de um piso
mínimo (\texttt{AUTOML\_MIN\_DETECTION\_RATE}). O \texttt{hit\_rate}
(equivalente ao \texttt{detection\_rate}) considera um alarme "detectado" se
existe ao menos um ponto anômalo dentro de uma janela de $\pm 1440$\,min
(24h) ao redor do horário do alarme -- \textbf{essa janela é simétrica},
o que será importante na Seção~\ref{sec:antecedencia}.

\subsection{Validação out-of-sample (OOS)}

Quando \texttt{AUTOML\_OOS\_SPLIT\_DATE} é definido, o ajuste do modelo
(treino, estatísticas de normalização e percentil de threshold) usa
\emph{apenas} dados anteriores a essa data; \texttt{hit\_rate},
\texttt{normal\_alert\_rate} e \texttt{composite\_score} são calculados
somente no período posterior -- alarmes e pontos que o modelo nunca viu.

% =====================================================================
\section{Histórico de experimentos}
\label{sec:historico}
% =====================================================================

A tabela~\ref{tab:historico} resume a sequência de execuções que levou ao
modelo final. Os números de \texttt{n\_alarms} nesta tabela refletem o que
a pipeline reportou originalmente (\emph{não} deduplicado -- ver ressalva
na Seção~\ref{sec:dados}).

\begin{table}[H]
\centering
\small
\begin{tabular}{@{}p{2.1cm}p{2.6cm}p{2.3cm}p{1.6cm}p{1.6cm}p{2.3cm}@{}}
\toprule
\textbf{Etapa} & \textbf{Config.} & \textbf{Melhor modelo} & \textbf{Hit rate} & \textbf{Composite} & \textbf{Observação} \\
\midrule
EXP5 (baseline) & split OOS 2025-05-01, grupo combinado & iforest & 57,1\% (40/70) & 0,857 & Ponto de partida \\
EXP5-v2 & split empurrado p/ 2025-07-01 & -- & -- & -- & \textbf{Achado negativo}: 92\% dos alarmes ficam antes de julho; \texttt{n\_alarms} colapsou p/ 8/4, scores viraram não comparáveis \\
EXP5-v3 & sem split OOS (in-sample), 3 grupos, dense[256,128] & dense & 78,3\%/81,6\% & 0,927/0,938 & Threshold do dense caiu p/ $\sim10^{-9}$ -- sinal de memorização, não generalização \\
EXP5-v4 & OOS restaurado (2025-05-01), mesmos 3 grupos & \textbf{iforest} (TC382) & \textbf{65,0\% (26/40)} & \textbf{0,883} & Vantagem do dense \emph{evaporou} OOS; TC382+iforest foi o único candidato estável nos dois regimes \\
Seed-sweep & TC382\_03\_A + iforest, 5 seeds & iforest & 65,0\% (todas) & -- & Variância de semente \textbf{zero} em \texttt{hit\_rate} \\
\bottomrule
\end{tabular}
\caption{Histórico resumido dos experimentos EXP5. Ver \texttt{docs/analise\_automl\_exp5.md} para o detalhamento completo de cada etapa.}
\label{tab:historico}
\end{table}

\subsection{Por que o dense com capacidade maior foi descartado}

Na tentativa de reproduzir o resultado da Lara (autoencoder denso
\texttt{[256, 128]}, ver \texttt{analise\_automl\_lara.md}), o \texttt{dense}
venceu com folga quando avaliado \emph{in-sample} (sem holdout OOS): até
78--82\% de \texttt{hit\_rate}. Porém o \texttt{threshold} vencedor ficou na
ordem de $10^{-9}$ -- praticamente zero -- compatível com um autoencoder de
grande capacidade memorizando quase perfeitamente um vetor de 1--2 features
contínuas sobre $\sim$720 mil pontos de treino, em vez de aprender um padrão
generalizável. Ao restaurar o split OOS, o \texttt{dense} e o \texttt{iforest}
empataram tecnicamente no grupo combinado (mesmo \texttt{hit\_rate}), e no
grupo \texttt{TC382\_03\_A} isolado o \texttt{iforest} superou o
\texttt{dense} de forma consistente -- confirmando que o ganho do dense era,
em grande parte, artefato de \emph{overfitting}.

% =====================================================================
\section{Modelo final selecionado}
\label{sec:modelo}
% =====================================================================

\begin{table}[H]
\centering
\begin{tabular}{@{}ll@{}}
\toprule
\textbf{Parâmetro} & \textbf{Valor} \\
\midrule
Grupo / sensor & \texttt{TC382\_03\_A\_univariado} (só \texttt{TC382\_03\_A}) \\
Modelo & \texttt{IsolationForest} (scikit-learn) \\
\texttt{contamination} & 0,05 \\
\texttt{n\_estimators} & 100 \\
\texttt{random\_state} & 42 (estável entre seeds, ver \S\ref{sec:seed}) \\
\texttt{threshold\_percentile} & 95,0 \\
\texttt{threshold} (erro de reconstrução) & 0,656 \\
\texttt{debounce} (pontos consecutivos) & 1 \\
Split OOS & 2025-05-01 (treino: jan--abr; avaliação: mai--out) \\
Máscara operacional & habilitada (\texttt{NGP\_A} como referência) \\
Normalização & \texttt{zscore}, estatísticas só do período de treino \\
\bottomrule
\end{tabular}
\caption{Configuração completa do modelo final.}
\end{table}

\subsection{Resultado de calibração (OOS, período mai--out/2025)}

\begin{table}[H]
\centering
\begin{tabular}{@{}lc@{}}
\toprule
\textbf{Métrica} & \textbf{Valor} \\
\midrule
\texttt{hit\_rate} & 65,0\% (26/40 -- valor da pipeline, não deduplicado) \\
\texttt{composite\_score} & 0,883 \\
\texttt{normal\_alert\_rate} & 2,43\% \\
\texttt{anomaly\_rate\_points\_per\_day} & 109,2 pontos/dia \\
\bottomrule
\end{tabular}
\end{table}

\subsection{Variância de semente}
\label{sec:seed}

Foi implementada uma checagem de variância (\texttt{AUTOML\_SEED\_SWEEP\_N},
adicionada em \texttt{automl\_pipeline.py} especificamente para este
experimento): o mesmo \texttt{iforest} vencedor foi re-treinado com 4 seeds
extras, mantendo \texttt{threshold\_percentile}/\texttt{debounce} fixos.

\begin{table}[H]
\centering
\begin{tabular}{@{}cccc@{}}
\toprule
\textbf{Seed} & \textbf{hit\_rate} & \textbf{normal\_alert\_rate} & \\
\midrule
42 (original) & 65,0\% & 2,43\% & \\
43 & 65,0\% & 2,47\% & \\
44 & 65,0\% & 2,03\% & \\
45 & 65,0\% & 1,90\% & \\
46 & 65,0\% & 1,86\% & \\
\midrule
\textbf{média / desvio} & \textbf{65,0\% / std=0} & 2,14\% / std=0,26pp & \\
\bottomrule
\end{tabular}
\caption{Variância de semente do candidato final. \texttt{hit\_rate} idêntico
(26/40) nos 5 seeds -- resultado estável, diferente do $\sim\pm27$pp de
ruído de semente observado no dense da pipeline da Lara.}
\end{table}

% =====================================================================
\section{Antecedência de detecção e falsos positivos}
\label{sec:antecedencia}
% =====================================================================

A métrica \texttt{hit\_rate} original usa uma janela \emph{simétrica} de
$\pm$24h ao redor do alarme -- ou seja, conta como "detectado" tanto uma
anomalia que precedeu o alarme (previsão real) quanto uma que só apareceu
\emph{depois} (o modelo reagiu ao mesmo evento, sem antecipar nada). Esta
seção separa os dois casos usando os alarmes já deduplicados
(20 alarmes reais no período OOS, ao invés de 40).

\subsection{Detecção: previsão vs.\ reação}

\begin{table}[H]
\centering
\begin{tabular}{@{}lcc@{}}
\toprule
\textbf{Categoria} & \textbf{Qtd.} & \textbf{\% dos 20 alarmes} \\
\midrule
Alarmes avaliados (OOS, deduplicados) & 20 & 100\% \\
Detectados dentro de $\pm$24h & 13 & 65,0\% \\
\quad -- com \textbf{antecedência real} (anomalia antes do alarme) & \textbf{9} & \textbf{45,0\%} \\
\quad -- só depois/simultâneo (reação, não previsão) & 4 & 20,0\% \\
Sem detecção & 7 & 35,0\% \\
\bottomrule
\end{tabular}
\caption{O \texttt{hit\_rate} de 65\% reportado pela pipeline mistura
previsão e reação. A taxa de previsão de fato é 45\% (9/20).}
\end{table}

\subsection{Tempo de antecedência (dos 9 casos preditivos)}

\begin{table}[H]
\centering
\begin{tabular}{@{}lc@{}}
\toprule
\textbf{Estatística} & \textbf{Valor} \\
\midrule
Mediana & 13,4h antes do alarme \\
Média & 15,9h \\
Mínimo & 0,7 min (praticamente simultâneo) \\
Máximo observado & 24,0h (censurado -- ver nota) \\
\bottomrule
\end{tabular}
\end{table}

\textbf{Nota sobre censura:} 4 dos 9 casos preditivos bateram exatamente no
limite da janela de avaliação (23,99--24,00h). Como a janela de matching é
fixa em $\pm$24h, a antecedência real desses casos pode ser
\emph{maior} que 24h -- não é possível medir além do limite da janela com a
metodologia atual.

\begin{figure}[H]
\centering
\includegraphics[width=0.85\textwidth]{pred_01_20250526_2129.png}
\caption{Exemplo de detecção com antecedência de $\sim$13,3h antes do alarme
de 26/05/2025. Linha verde tracejada = horário do alarme; linha vermelha
pontilhada = primeira anomalia detectada; pontos vermelhos = todos os pontos
sinalizados como anômalos na janela.}
\end{figure}

\begin{figure}[H]
\centering
\includegraphics[width=0.85\textwidth]{pred_05_20250531_0612.png}
\caption{Exemplo de detecção com antecedência de $\sim$10,2h antes do alarme
de 31/05/2025.}
\end{figure}

\begin{figure}[H]
\centering
\includegraphics[width=0.85\textwidth]{pred_09_20250617_1649.png}
\caption{Caso-limite: antecedência de apenas 0,7 min (43 segundos) --
detecção praticamente simultânea ao alarme de 17/06/2025, não uma
previsão real.}
\end{figure}

\subsection{Falsos positivos}
\label{sec:fp}

Falso positivo aqui é definido como ponto anômalo em operação
(\texttt{operational\_state == "on"}), fora de qualquer janela de
$\pm$24h ao redor de \emph{qualquer} alarme real de \texttt{TC382\_03\_A}
(considerando os 38 alarmes reais do ano todo, não só os 20 do período OOS).
Pontos anômalos contíguos (gap $\le$ 5min) foram agrupados em episódios.

\begin{table}[H]
\centering
\begin{tabular}{@{}lc@{}}
\toprule
\textbf{Métrica} & \textbf{Valor} \\
\midrule
Episódios de falso positivo (período OOS, $\sim$183 dias) & 63 \\
Taxa & $\approx$0,34 episódios/dia (1 a cada $\sim$3 dias) \\
Duração mediana do episódio & 2,5 min \\
Duração no 75º percentil & 16 min \\
Duração máxima observada & 2416 min ($\sim$40,3h) \\
\bottomrule
\end{tabular}
\end{table}

A distribuição de duração é fortemente assimétrica: a maioria dos episódios
é um pico isolado de poucos minutos (ruído do sensor ou transição breve não
capturada pela máscara operacional), mas um pequeno número de episódios
sustentados por horas a dias concentra a maior parte dos pontos marcados
como falso positivo.

\begin{figure}[H]
\centering
\includegraphics[width=0.8\textwidth]{01_fp_duration_histogram.png}
\caption{Distribuição de duração dos 63 episódios de falso positivo (escala log).}
\end{figure}

\begin{figure}[H]
\centering
\includegraphics[width=0.85\textwidth]{fp_01_20250817_1740.png}
\caption{Maior episódio de falso positivo do período: $\sim$40h contínuas
entre 17 e 19/08/2025, sem alarme real correspondente registrado. Vale
investigar operacionalmente o que ocorreu nesse intervalo.}
\end{figure}

\begin{figure}[H]
\centering
\includegraphics[width=0.85\textwidth]{fp_03_20250527_2131.png}
\caption{Episódio de falso positivo de $\sim$9,5h em 27--28/05/2025.}
\end{figure}

\subsection{Visão geral do período}

\begin{figure}[H]
\centering
\includegraphics[width=\textwidth]{00_overview.png}
\caption{Série bruta de \texttt{TC382\_03\_A} no período OOS completo
(mai--out/2025). Sombra cinza = \texttt{off\_longo} (parada prolongada real);
traços laranja no rodapé = \texttt{off\_curto}/\texttt{transiente}
(episódios curtos de $\sim$1h, marcados como referência pois são curtos
demais para sombrear nessa escala de tempo sem distorcer a leitura visual);
linhas verdes tracejadas = alarmes reais; pontos vermelhos = anomalias
detectadas pelo modelo.}
\end{figure}

% =====================================================================
\section{Limitações}
% =====================================================================

\begin{itemize}
  \item \textbf{Amostra pequena de alarmes:} apenas 20 alarmes reais no
        período de avaliação OOS. Frações como 65\%, 45\% têm incerteza
        estatística considerável ($\pm$1 alarme já move o resultado em 5pp).
  \item \textbf{Janela de matching simétrica:} a métrica original
        (\texttt{hit\_rate}) não distingue previsão de reação -- corrigido
        nesta análise, mas a pipeline de produção ainda reporta o número
        combinado por padrão.
  \item \textbf{Censura em 24h:} antecedências acima de 24h não são
        mensuráveis com a janela atual; quase metade dos casos preditivos
        bateu nesse teto.
  \item \textbf{Deduplicação de alarmes:} feita manualmente nesta análise;
        ainda não incorporada ao carregamento de dados da pipeline
        (\texttt{src/cnn1d\_ae/io.py}) -- os \texttt{n\_alarms} reportados
        em \texttt{calibration\_report.json} continuam em dobro do real até
        essa correção ser aplicada.
  \item \textbf{Univariado:} o modelo usa somente \texttt{TC382\_03\_A};
        alarmes de outros sensores do equipamento estão fora do escopo
        desta análise.
\end{itemize}

% =====================================================================
\section{Conclusão e próximos passos}
% =====================================================================

O candidato \texttt{TC382\_03\_A} + \texttt{IsolationForest}, com a
configuração da Tabela~2, é até o momento o único resultado desta série de
experimentos validado tanto in-sample quanto out-of-sample, robusto contra
as alternativas testadas (\texttt{dense}, \texttt{ocsvm}) e estável entre
sementes aleatórias. Sua taxa de detecção real com antecedência (45\%,
mediana de 13,4h) é o número mais relevante para uso operacional -- não o
\texttt{hit\_rate} agregado de 65\%, que inclui reações tardias.

\textbf{Próximos passos sugeridos:}
\begin{enumerate}
  \item Corrigir a deduplicação de alarmes na pipeline (\texttt{io.py}),
        para que \texttt{n\_alarms} reflita a contagem real por padrão.
  \item Investigar operacionalmente o episódio de $\sim$40h de falso alerta
        (17--19/08/2025) -- não corresponde a nenhum alarme real registrado.
  \item Repetir a análise de antecedência/falsos-positivos com janela de
        matching assimétrica nativa na pipeline (só "antes"), em vez de
        pós-processamento manual como feito aqui.
  \item Avaliar se o mesmo candidato (\texttt{iforest} univariado) se
        sustenta em outros sensores/equipamentos do Cabiúnas.
\end{enumerate}

\end{document}
